
Pour le premier il s'agit de transformer \(s\) en la matrice unité \(\underline{\underline{1}}=(\delta_{ij})\) dans un voisinage zariskien de \(y\) en multipliant \(s\) à gauche et à droite par des matrices convenables \(c_{h,k,\lambda}\). Montrons d'abord qu'on peut rendre ainsi \(s_{11}=1\). Si \(s_{11}\ne1\) et \(s_{21}\in\underline{p}\), nous cherchons une matrice \(c=c_{h,1,\lambda}\) tel que le produit \(s'=sc=(s'_{ij})\) ait \(s'_{21}\notin\underline{p}\). Puisque \(\det(s)=1\notin\underline{p}\), il existe au moins un indice \(i\) tel que \(s_{i1}\notin\underline{p}\) Alors on a \(s_{21}+s_{i1}\notin\underline{p}\), et \(c_{i,1,1}\) répond à la question. On peut donc supposer \(s_{21}\notin\underline{p}\), et quitte à localiser, on peut supposer \(s_{21}\) inversible. Si \(s_{11}-1\in\underline{p}\), alors dans \(s'=(s'_{ij})=sc_{1,2,1}\) on a
\(s'_{11}-1=(s_{11}+s_{21})-1\notin\underline{p}\). Quitte à localiser, on peut alors supposer \(s_{11}-1\) inversible. Alors \(\lambda=(1-s_{11})/s_{21}\) est inversible et on trouve \(s'_{11}=1\) dans \(s'=(s'_{ij})=sc_{1,2,\lambda}\).

Supposant \(s_{11}=1\) on peut alors rendre \(s\) à la forme où
\[
 s_{11}=1,\qquad s_{1i}=s_{i1}=0\quad(i=2,\ldots,n)
\]
dans un voisinage zariskien, de \(y\) en multipliant \(s\) à gauche (resp. à droite) par des matrices \(c_{i,1,\lambda}\) (resp. \(c_{i,1,\lambda}\)). Puis on va rendre \(s_{22}=1\) etc. Par récurrence on peut alors transformer \(s\) dans un voisinage zariskien \(U\) de \(y\) en une matrice \(s'=(s'_{ij})\) ayant \(s'_{ij}=\delta_{ij}\) pour \(i<n\) ou \(j<n\). Mais par construction même on a alors \(s'=c_1sc_2\) sur \(\mathcal U\), où \(c_1\) et \(c_2\) sont des produits de matrices \(c_{h,k,\lambda}\). Donc on a
\[
 \det(s')=\det(s)=1,
\]
de sorte que \(s'_{nn}=1\), \(s'=\underline{\underline{1}}\), et on a gagné.

On suppose donc que \(s\) soit déjà sur \(Y\) un produit de matrices \(c_{h,k,\lambda}\).
Reste à voir que celles-ci sont des produits de commutateurs dans un voisinage étale de \(y\). Tout d'abord on voit que deux matrices
